\subsection{Core Metadata and Validation Proofs}
\label{app:validation-proofs}
This appendix defines the inference rules used by the formal validator in
Section~\ref{sec:candidate-validation}. It proves soundness of this calculus;
it does not identify a diagnostic trace from the implementation with a
checked proof certificate.

\subsubsection{Finite-Sequence Interpretation}
A metadata valuation assigns a value of the declared sort to every captured
field and shape. Axis sequences contain identities with associated extents.
All sequences are finite. For $Q=[Q_0,\ldots,Q_{n-1}]$,
$\gath(Q,q)=[Q_{q_0},\ldots,Q_{q_{k-1}}]$ when every position is in range.
$\pos(Q,G)$ returns the unique positions of $G$'s members in $Q$, in $G$'s
order, and requires distinct $Q,G$ and inclusion of their member sets.
Ordered complement deletes members while preserving the first argument's
order. For distinct sequences, $\mpo(A,B)=(u,v)$ lists the positions of
common members in $A$'s order and their corresponding positions in $B$.
The inverse of a permutation $q$ is the unique list $q^{-1}$ satisfying
$q_{q^{-1}_i}=i$. Extents and ranks never depend on tensor entries or
reduction coordinates.

More explicitly, the typed core expression grammar is
\begin{equation}
\begin{aligned}
 e::={}&\operatorname{src}(k)\mid c\mid[e_1,\ldots,e_m]
       \mid e_1\Vert e_2\mid |e|\mid\operatorname{range}(e)\\
      &\mid\gath(e,e)\mid\ocomp(e,e)\mid\pos(e,e)
       \mid\operatorname{inverse}(e)\\
      &\mid\mpo(e,e)_i\mid e+e\mid e-e
       \mid\operatorname{if}(b,e,e).
\end{aligned}
\end{equation}
Here $\operatorname{src}(k)$ names captured metadata, $c$ is a typed
literal, and $i\in\{1,2\}$. Arithmetic is scalar or an integer shift of a
position sequence, as determined by the types. Predicates and their Boolean
combinations are the $b$ expressions. Sequence member-set operations in
the reasoning rules are ordered filter/union views, with the member-set
interpretations below; they introduce no new axis identities.

Metadata evaluation follows these definitions. Literals and captured values
are defined; lists, concatenation, and length are strict.
$\operatorname{range}(n)=[0,\ldots,n-1]$ requires $n\ge0$ and representable
length. Complement and paired matching require finite, distinct arguments;
permutation inverse requires $\operatorname{Perm}(q,|q|)$. Position shifts
and integer arithmetic additionally check representability in their declared
integer type. An inverse or position lookup with unmet premises is
undefined. Conditionals and Boolean short-circuiting evaluate only the
selected subexpression. A resource limit may interrupt an implementation;
the mathematical evaluator describes its result with sufficient resources.

\subsubsection{The Core Inference System}
We write $e\Downarrow_p v$ for successful evaluation of a metadata
expression. Evaluation is strict except for conditionals and short-circuit
Boolean operators; integer overflow is undefined.

Let $\Gamma$ contain legal LHS metadata and the fact that the written
condition evaluated to true. Derived facts must follow from these
premises. The judgments
\begin{equation}
 \Gamma\vdash e\downarrow,\qquad \Gamma\vdash\phi
 \label{eq:metadata-judgments}
\end{equation}
establish defined evaluation and a true metadata predicate, respectively.
These are derivability judgments; $e\Downarrow_p v$ describes evaluation
at one metadata valuation, not a proof-search outcome.
Predicates include ordered equality, length and integer comparisons,
distinctness, set inclusion, disjointness, and selector validity.
Set predicates forget order; ordered equality does not.

For a sequence $Q$, $\operatorname{Sel}(q,Q)$ requires distinct positions
in $[0,|Q|)$; $\operatorname{Perm}(q,n)$ also requires $|q|=n$.
Table~\ref{tab:metadata-definedness} gives rules for well-typed expressions,
with definedness premises explicit. Injectivity is a core restriction;
ordinary gather permits repeated positions.

\begin{table}[!htbp]
\caption{Representative inference rules for candidate validation.}
\label{tab:metadata-definedness}
\small
\begin{tabularx}{\linewidth}{@{}lY@{}}\toprule
Rule & Derivation\\\midrule
\textsc{D-Gather} &
$\displaystyle
 \frac{\Gamma\vdash Q\downarrow\quad\Gamma\vdash q\downarrow
       \quad\Gamma\vdash\operatorname{Sel}(q,Q)}
      {\Gamma\vdash\gath(Q,q)\downarrow}$\\[6pt]
\textsc{D-If-True} &
$\displaystyle
 \frac{\Gamma\vdash b\quad\Gamma\vdash e_1\downarrow}
      {\Gamma\vdash\operatorname{if}(b,e_1,e_2)\downarrow}$\\[6pt]
\textsc{E-Cong} &
$\displaystyle
 \frac{\Gamma\vdash e=e'\quad\Gamma\vdash H(e)\downarrow
       \quad\Gamma\vdash H(e')\downarrow}
      {\Gamma\vdash H(e)=H(e')}$\\[2pt]\bottomrule
\end{tabularx}
\end{table}
Here $H$ is a core operator with other arguments fixed; the false-branch
rule is symmetric. If $\Gamma$ proves $Q,q,h$ defined and
$\operatorname{Sel}(q,Q),\operatorname{Sel}(h,q)$, gather composition gives
\begin{equation}
 \Gamma\vdash\gath(\gath(Q,q),h)=\gath(Q,\gath(q,h))
 \quad\text{\textsc{E-Gather}}.
 \label{eq:validated-gather}
\end{equation}
Both sides have entry $Q[q[h[j]]]$ at every finite length.
The remaining rules and their soundness proofs follow below.

Definedness is checked on the original constructor expressions before
value simplification can erase a partial subexpression. For example,
eager evaluation of $0\cdot q[i]$ still requires $i$ to be in range.
Merely obtaining a simplified constant zero is insufficient. Supplied
fields are checked and preserved just like inferred fields.


The context $\Gamma$ initially contains legal source metadata and the
successfully true written guard. Source legality contributes the declared
operation contracts, including in-range selectors, permutations, distinct
input frames, paired lengths and extents, and output-layout equations.
These are assumptions about legal source instances. The same facts for an
RHS operation must be derived before they can be used.

The following rules define the core; no rule may assume an unproved
candidate obligation.
\begin{enumerate}
\item \emph{Premises and Boolean reasoning.} Use a premise from $\Gamma$,
introduce or eliminate conjunctions, and apply propositional consequences
validated by a finite Boolean truth table over defined atomic predicates.
A guard may contribute facts only along evaluation paths established by
its successful result; a short-circuited branch contributes no facts or
definedness. In particular, a successfully true conjunction establishes
both conjuncts, whereas a true disjunction does not establish either
disjunct separately. A premise involving a partial expression includes
successful evaluation of that expression.
\item \emph{Definedness.} Use Table~\ref{tab:metadata-definedness} and the
finite-sequence contracts above. Literals, source values, lists, length,
and concatenation compose strictly. Position lookup additionally uses
\begin{equation*}
 \frac{\begin{gathered}
       \Gamma\vdash Q\downarrow\quad\Gamma\vdash G\downarrow\\
       \Gamma\vdash\operatorname{distinct}(Q)\land
                    \operatorname{distinct}(G)\land
                    \operatorname{set}(G)\subseteq\operatorname{set}(Q)
       \end{gathered}}
      {\Gamma\vdash\pos(Q,G)\downarrow}
 \quad\text{\textsc{D-Pos}}.
\end{equation*}
The false-branch counterpart of \textsc{D-If-True} uses
$\Gamma\vdash\neg b$ and $\Gamma\vdash e_2\downarrow$.
If $\Gamma\vdash b\downarrow$ but its truth value is unresolved, requiring
both branches defined suffices; the same approach applies to short-circuit
operators.
\item \emph{Equality.} Reflexivity applies to defined expressions.
Symmetry, transitivity, substitution, and \textsc{E-Cong} require all
involved expressions to be defined; congruence applies argumentwise.
Ground operations use the finite-sequence definitions directly.
\item \emph{Sequence laws.} Instantiate the laws in
Table~\ref{tab:core-sequence-laws}, with their stated premises and
argument definedness. $\operatorname{set}(Q)$ denotes the member set.
\item \emph{Set reasoning.} A sequence intersection, difference, or union
has the corresponding member-set interpretation; its order is inherited
from its specified traversal. To prove inclusion, disjointness, or set
equality, replace each atomic member set by a Boolean variable for an
arbitrary member. Accept only a propositional consequence of the
established set premises. This finite truth-table procedure proves member
relations, never ordered equality.
\item \emph{Integer bounds.} Use exact ground computation and interval
bounds. If $a\in[l_a,u_a]$ and $b\in[l_b,u_b]$, then
$a+b\in[l_a+l_b,u_a+u_b]$ and
$a-b\in[l_a-u_b,u_a-l_b]$. Lengths are nonnegative;
concatenation lengths add. A computed interval proves representability
only when it lies inside the declared integer range. Unsupported arithmetic
or an overflow while computing a bound gives no derivation.
\end{enumerate}
These rules specify a proof relation, not a complete strategy for finding
proofs. Finite work bounds can abandon search without changing that relation.

\begin{table}[!htbp]
\caption{Sequence laws of the validation core. All expressions must be defined.}
\label{tab:core-sequence-laws}
\small
\begin{tabularx}{\linewidth}{@{}p{0.56\linewidth}Y@{}}\toprule
Law & Further premises\\\midrule
$|A\Vert B|=|A|+|B|$; $|\gath(A,q)|=|q|$
& Valid gather selector\\
$\gath(A,\operatorname{range}(|A|))=A$
& None\\
Gather composition, Equation~\eqref{eq:validated-gather}
& $\operatorname{Sel}(q,A)$ and $\operatorname{Sel}(h,q)$\\
$\gath(A,q\Vert h)=\gath(A,q)\Vert\gath(A,h)$
& All listed selectors in range\\
$\gath(A\Vert B,q)=\gath(A,q)$
& $\operatorname{Sel}(q,A)$\\
$\gath(A\Vert B,|A|+q)=\gath(B,q)$
& $\operatorname{Sel}(q,B)$\\
$\gath(c+\operatorname{range}(n),q)=c+q$
& $q$ selects positions in $\operatorname{range}(n)$\\
$\gath(A,\pos(A,G))=G$
& $A,G$ distinct and $\operatorname{set}(G)\subseteq\operatorname{set}(A)$\\
$\gath(q,\operatorname{inverse}(q))=\operatorname{range}(|q|)$
& $\operatorname{Perm}(q,|q|)$\\
$\operatorname{Perm}(\gath(q,h),n)$
& $\operatorname{Perm}(q,n)$ and $\operatorname{Perm}(h,n)$\\
$\operatorname{set}(\gath(A,q))\subseteq\operatorname{set}(A)$;
$\operatorname{distinct}(\gath(A,q))$
& $A$ distinct and $\operatorname{Sel}(q,A)$\\\bottomrule
\end{tabularx}
\end{table}

The finite matching and complement operations also publish their defining
properties: for $(u,v)=\mpo(A,B)$, the selectors are valid and
$\gath(A,u)=\gath(B,v)$ lists exactly the common members; a selector and
its ordered complement are disjoint and cover the original positions.
These facts follow from their definitions, not from the candidate rewrite.
We also use the finite length bounds implied by selectors and complements:
an injective selector into $Q$ has length at most $|Q|$, and deleting
positions cannot increase a sequence's length. Nonnegative block lengths
bound each block offset and position by the total length. These inequalities
follow by injection into the original position set and monotonicity of
addition; a proved bound on the total therefore bounds each such offset.

\begin{lemma}[Sound metadata derivations]
\label{lem:metadata-soundness}
For every valuation $p$ satisfying $\Gamma$, a derivation
$\Gamma\vdash e\downarrow$ implies successful evaluation of $e$ at $p$.
A derivation $\Gamma\vdash\phi$ implies successful evaluation of $\phi$
to true at $p$.
\end{lemma}
\begin{proof}
Induct on the derivation. Premises hold by satisfaction of $\Gamma$;
propositional rules preserve truth. For definedness, a valid gather selector
makes every accessed index in range; distinctness and inclusion make every
position lookup unique and present; a permutation has a unique inverse.
All other sequence constructors process finitely many defined arguments.
The integer-interval rules follow from monotonicity of addition and
subtraction and establish both exact values and representability.
Conditional rules check either both branches or the established selected
branch, so actual evaluation succeeds.

Equality rules preserve equality of already defined values. The sequence
laws follow by comparing lengths and entries: gather composition has entry
$A_{q_{h_j}}$ on both sides; concatenation and its shifted selectors agree
within each block; position lookup selects exactly the requested members;
inverse selection gives the identity list. Composition of two permutations
is a bijection. Injective selection preserves distinctness and membership.
Selecting a shifted range gives entry $c+q_j$ on both sides.
Matching visits precisely the common members and records their paired
locations; complement visits precisely the omitted positions. Finally,
for any member, the Boolean interpretation of set operations is exact.
A propositional consequence therefore holds for that member, and hence
for every member. None of these cases fixes a sequence length.
\end{proof}

\begin{lemma}[Selection and complement transport]
\label{lem:selection-complement}
Let $q,r$ be disjoint injective selectors into a sequence $A$, and set
$s=\ocomp(\operatorname{range}(|A|),r)$. Then $\pos(s,q)$ is defined,
\begin{equation}
 \gath(\gath(A,s),\pos(s,q))=\gath(A,q),
 \label{eq:complement-selection}
\end{equation}
and deleting those selected entries from $\gath(A,s)$ gives the same
ordered sequence as deleting $r\Vert q$ from the positions of $A$.
\end{lemma}
\begin{proof}
Disjointness places every entry of $q$ in $s$; both lists are distinct.
The position-lookup rule proves definedness, and gather composition with
$\gath(s,\pos(s,q))=q$ proves Equation~\eqref{eq:complement-selection}.
For the deletion claim, each original position remains exactly when it is
in neither $r$ nor $q$. Both constructions visit these positions in their
original increasing order, so they have the same entries in the same order.
This argument concerns positions and does not require distinct values in
$A$. It also applies to empty selections. We use this finite-sequence
consequence as a derived transport rule in the reassociation proof.
\end{proof}

\subsubsection{Validation Procedure}
\label{app:validation-procedure}
Let $\mathcal G$ be the finite set of candidate obligations. These are
proof goals, never source assumptions. The procedure is:

\noindent\begin{minipage}{\linewidth}
\begin{lstlisting}[language={},caption={Validation of a completed candidate.}]
Validate(L, R, C, f, goals):
    Gamma = source legality and successful evaluation of C
    require source-only, typed, defined original expressions in f
    require supplied fields are preserved and all holes are filled
    prove every goal in goals from Gamma
    lower L and completed R using the declared operation contracts
    prove every RHS legality and coordinate-transport obligation
    normalize both value expressions using checked algebraic laws
    prove equality of normalized terms, ordered outputs and shapes
    return Proved if all goals have derivations; otherwise Unknown
\end{lstlisting}
\end{minipage}


Here \code{goals} denotes $\mathcal G$. A \code{Proved} result requires
all obligations to have derivations; unresolved goals or work limits give
\code{Unknown}. This outcome does not establish infeasibility.

For the canonical program of Theorem~\ref{thm:symbolic-compilation},
acceptance instead requires guard certification:
\begin{equation}
 \forall p.\quad \operatorname{Legal}_L(p)\land C(p)
 \Rightarrow\operatorname{check}(p).
 \label{eq:guard-feasibility}
\end{equation}
This obligation uses that theorem's domain assumptions. Candidate
validation applies to any completed program with a derivation in the
calculus. Implementation paths outside the calculus are discussed in
Section~\ref{sec:implementation-scope}.

\subsubsection{Mappings, Partitions, and Executable Expressions}
\begin{lemma}[Mapping and partition preservation]
\label{lem:validated-mappings}
Mappings derived by valid selection, gather composition, concatenation
with offsets, permutation reconstruction, and paired overlap matching
satisfy their stated ordered-selection equations. A registered complete
partition covers its interface exactly once.
\end{lemma}
\begin{proof}
For selection and composition, use Table~\ref{tab:core-sequence-laws}.
For reconstruction, concatenate groups as $\vec G=\vec G_1\Vert\cdots\Vert \vec G_m$
and their position lists as $q=q_1\Vert\cdots\Vert q_m$. Require $q$ to
be a permutation and $|\vec G_i|=|q_i|$. With
$Q=\gath(\vec G,\operatorname{inverse}(q))$, inverse composition gives
$\gath(Q,q)=\vec G$, whose corresponding blocks give
$\gath(Q,q_i)=\vec G_i$. Injectivity and completeness of $q$ give exact coverage.
For overlap, if $\vec G=\gath(Q,q)$, $\vec H=\gath(Q,h)$, and
$(u,v)=\mpo(q,h)$, then $q_{u_j}=h_{v_j}$; both transported selections
therefore have entry $Q_{q_{u_j}}$. The paired selectors retain order and
partner identity. Induction on a mapping derivation proves the claim.
\end{proof}

\begin{lemma}[Core metadata emission]
\label{lem:metadata-emission}
Translate a defined core metadata expression to executable metadata syntax
by preserving source references, literals, operator identities, argument
order, and conditional evaluation order. With the primitive sequence and
checked-integer semantics above, executable evaluation has the same value
as metadata evaluation.
\end{lemma}
\begin{proof}
Induct on the expression. Source references read the same captured metadata;
literals are unchanged. For a strict operator, the induction hypotheses
give equal argument values. The corresponding executable operator applies
the same finite-sequence definition or checked arithmetic operation.
Definedness rules ensure that its accesses and representability checks
succeed. For a conditional or short-circuit operator, both evaluators first
obtain the same Boolean and evaluate the same selected subexpression,
whose value agrees by induction. The tree is finite and each primitive
terminates on finite arguments. Thus execution agrees, with sufficient
resources. This proves the specified translation, not correctness of every
native or generated backend implementation of a primitive.
\end{proof}

\subsubsection{Lowering, Normalization, and Comparison}
A lowered value expression uses constants, indexed input reads, addition,
multiplication, and finite summation. Its normalized form is a sum of
terms $(c,[r_1,\ldots,r_m],\vec D)$, interpreted by
Equation~\eqref{eq:validator-term}.
\begin{equation}
 \sum_{\rho\in\operatorname{Coords}(\vec D)}c\prod_{i=1}^{m}r_i(\rho).
 \label{eq:validator-term}
\end{equation}
The read list preserves multiplicity. Normalization uses distributivity
and reduction linearity, combining reductions only with proved disjoint
bound axes and independence of moved factors. Bound-coordinate renaming
must be injective, preserve extents, and avoid capture. Coefficient
computation is exact; representation overflow leaves the goal unresolved.

Term comparison checks equal coefficients, matching read multisets in a
common coordinate frame, and the same reduction axes and domains. Output
comparison additionally establishes ordered axis equality, shape equality,
and equality of the source and replacement result dtypes. Equality of shapes alone does not establish equality
of reads or axis order. The normalization and comparison proof follows below.


\begin{lemma}[Core lowering]
\label{lem:core-lowering}
Lowering a legal core expression using validated coordinate mappings gives
its declared ordered shape and its value at every output coordinate.
\end{lemma}
\begin{proof}
Induct on the expression. An input read uses its physical dimension order.
For a permutation $q$, output coordinate $i$ supplies the coordinate of
input dimension $q_i$; the permutation premise makes this a bijective
reindexing. Unary summation binds precisely the selected dimensions and
preserves the remaining dimension order. A contraction pairs its batch and
contraction axes in attribute-list order. Equal paired extents provide a
common domain, distinct disjoint selectors avoid duplicate bindings, and
the declared batch/left-free/right-free order specifies the output. Its
value is the finite sum of the product of the corresponding operand
entries. Addition requires matching ordered shapes and adds corresponding
entries. Substitute lowered operands using the induction hypotheses;
Lemma~\ref{lem:validated-mappings} preserves their coordinates. RHS
legality facts must already have derivations, so no illegal replacement is
used as a premise.
\end{proof}

Define $\operatorname{deps}(t)$ as a conservative set of axes occurring in
a term's read coordinates; absence implies independence. Bound names are
fresh, and any renaming is a capture-avoiding, extent-preserving bijection.
The core normalizer maps a constant to one coefficient term, a read to one
read term, addition to term-list concatenation, and reduction to prepending
its bound axes after a disjointness check. A product forms the Cartesian
product of its child term lists. For terms with reduction axes $\vec D_1,\vec D_2$
and read factors $F_1,F_2$, this step requires
\begin{equation}
\begin{aligned}
 \vec D_1\cap \vec D_2&=\varnothing,\\
 \vec D_1\cap(\operatorname{deps}(F_2)\setminus \vec D_2)&=\varnothing,\\
 \vec D_2\cap(\operatorname{deps}(F_1)\setminus \vec D_1)&=\varnothing.
\end{aligned}
\end{equation}
Intersections here refer to member sets. All three conditions must be
proved; normalization does not assume independence from a missing map.

\begin{lemma}[Core normalization and comparison]
\label{lem:core-normalization}
Successful core normalization preserves values. Matching normalized terms
with equal coefficients, equal read multisets, and equal reduction domains
establishes value equality.
\end{lemma}
\begin{proof}
Induct on the normalized value expression. Constants and reads are immediate;
addition is concatenation of finite sums. Reduction distributes over a
finite sum. Disjointness permits nested reduction domains to be identified
with their Cartesian product. For multiplication, the stated independence
conditions justify
\begin{equation}
 \left(\sum_{\vec D_1}a\right)\left(\sum_{\vec D_2}b\right)
 =\sum_{\vec D_1\Vert \vec D_2}ab.
\end{equation}
Distributivity gives the Cartesian product of term lists; exact coefficient
arithmetic and multiplication preserve each term. These identities hold
also for empty domains. Collecting equal terms adds coefficients in exact
real arithmetic. A legal renaming is a bijection between coordinate
assignments, preserving every summand. For matching terms, multiplication
is commutative, read multiplicities agree, and summation order does not
change a finite sum. Therefore each matched pair has equal value, as do
the complete term lists. This is sufficient comparison, not a completeness
claim about normal forms or proof search.
\end{proof}

\begin{proof}[Proof of Theorem~\ref{thm:validator-soundness}]
Fix legal source metadata $p$ for which the written guard $C$ evaluates to
true. A \code{Proved} result means that every required check has a
derivation. By Lemma~\ref{lem:metadata-soundness}, the attribute expressions
are defined, supplied fields are preserved, and every RHS operation is
legal. Lemma~\ref{lem:metadata-emission} shows that executing $f$ therefore
produces the established attributes $\beta$, with sufficient resources.

Lemma~\ref{lem:core-lowering} gives each lowered graph its declared shape
and value at every output coordinate. By
Lemma~\ref{lem:core-normalization}, normalization preserves these values
under exact arithmetic, and successful term comparison establishes equal
values for all compatible tensor inputs. The proved ordered-shape checks
and result-dtype checks give the remaining parts of
$\operatorname{Valid}(p,\beta)$. Thus
$f\Downarrow_p\beta\land\operatorname{Valid}(p,\beta)$.
Since $p$ was arbitrary, the theorem follows independently of how the
candidate was generated.
\end{proof}
A new implementation rule must be justified by this calculus or receive
its own soundness argument before it belongs to this formal guarantee.

\subsubsection{Validation of Transpose Composition}
\label{app:transpose-validation}
For transpose composition on input $t$, let $a,b$ be source permutations,
$n=|\operatorname{Ord}_t|$, and $f(p)=\gath(a,b)$. Source legality and
\textsc{D-Gather} prove definedness; permutation composition gives the left
premise below. \textsc{E-Gather} gives the right, with
$F_L=\gath(\gath(\operatorname{Ord}_t,a),b)$ and
$F_R=\gath(\operatorname{Ord}_t,f(p))$.
\begin{equation}
 \frac{\Gamma\vdash\operatorname{Perm}(f(p),n)
       \qquad\Gamma\vdash F_L=F_R}
      {\Gamma\vdash\operatorname{Valid}(p,f(p))}
 \quad\text{\textsc{Transpose}}
 \label{eq:transpose-validation-tree}
\end{equation}
This derived rule is specific to the transpose-only rewrite: both sides
read the same input in the common frame and preserve its dtype. No
reductions occur, and the argument includes rank zero.

\subsubsection{Validating Batched Reassociation}
\label{app:reassociation-validation}
We validate the complete program in
Table~\ref{tab:reassociation-program}, rather than treating the candidate's
axis correspondences as premises. Let $C$ be the batch-preservation and
integer-bound guard in Appendix~\ref{app:batched-reassociation-program}, and let
$\Gamma$ contain only $C$ and source legality. The argument uses the core
sequence laws, selection/complement transport
(Lemma~\ref{lem:selection-complement}), and finite-sum normalization.

\begin{proposition}[Guarded reassociation program]
\label{prop:reassociation-program}
For every legal source match satisfying $C$, the program in
Table~\ref{tab:reassociation-program} evaluates to legal attributes and
preserves the ordered output shape and values under exact real arithmetic.
Ranks, batch lengths, and contraction lengths may vary; free groups may
be interleaved in the operands. Zero and singleton extents are allowed.
\end{proposition}
\paragraph{Validation obligations.}
The proof discharges three concrete obligations from source legality and
$C$, without assuming candidate correspondences:
\begin{enumerate}
\item \emph{Defined attributes.} The partition places $z_I$ inside $f_Z$
and $k_Y$ inside $f_Y$, so both critical position queries are defined.
The integer bound makes their offsets representable.
\item \emph{Legal replacement.} Paired selectors are injective, in range,
and disjoint from batching; they select $\Jsel$ in the inner replacement
and $\Isel\Vert\vec K$ in the outer one, with matching extents.
\item \emph{Equivalent result.} Both outputs have ordered frame
$\vec B\Vert\Irem\Vert\Jrem\Vert\vec H$, where $\vec H$ is the outer
source dot's right-free group. An extent-preserving bijection splits the
source reduction group $\vec D$ into its selected $\vec I$ and $\vec J$
parts; checked finite-sum laws then give identical reads and reductions.
\end{enumerate}
Appendix~\ref{app:reassociation-proof} supplies the full coordinate and
normalization derivation. Appendix~\ref{app:reassociation-record} connects
the displayed program to recorded compiler output; the argument quantifies
over arbitrary legal metadata rather than sampled ranks.

\subsubsection{Suggesting Applicability Conditions}
\label{app:condition-suggestions}
When validation under $C$ is unresolved, TEPL can propose a stronger guard
$C\land G$. A bounded search tries individual predicates and pairs drawn
from captured-axis equalities, producer batch/free positions, and identity
permutation restrictions. Each trial reruns construction and validation in
a fresh context, proving correctness for every legal match under $C\land G$;
a separate concrete search must find a legal satisfying match to establish
nonvacuity. Failed validation or witness search yields no suggestion, and
removing a predicate requires both checks again. The user must add the
suggestion to \code{where} and revalidate. Neither weakest conditions nor
complete discovery are claimed. A worked reassociation suggestion is
archived in \nolinkurl{evaluation/supplementary-materials/}.

